% ch3_g 第3章 §3.6尾–§3.8 (书页 134–143 / p149–p158)
% 说明：本块自书页 134（p149）顶部起。p149 顶部为图 3.15 及一个完整段落，
% 并非从书页 133 延续来的半句，故不用 %JOIN% 标记。本块为章末块，书页 143
% 自然收束到第 3 章结尾（下一页 p159 为第 4 章章首），故不用 %--TAIL-- 标记。

% ---- 书页 134 (p149) ----
%FIG{3.15}: {(a) 凝聚体的相位绕圆环扭转 $n$ 次，代表一个超流流动。(b) 一个涡旋可以把相位扭转从 $n$ 次变为 $n-1$ 次。}

有趣的是，自由玻色子的 $v_{c}=0$，即使在零温下也是如此。因此，尽管零温下的自由玻色子具有长程序，但由于 $v_{c}=0$，它们并不形成超流体。

让我们更仔细地考察二流体图像（two-fluid picture）。注意 $m\pmb{v}$ 决定凝聚体的扭转。因此，超流分量中的每个玻色子都携带动量 $m\pmb{v}$，而 $\pmb{P}_{\pmb{v}}/m\pmb{v}$ 给出超流分量中玻色子的数目。我们求得超流密度（superfluid density）为
\begin{equation*}
\rho_{s} = \frac{\pmb{P}_{\pmb{v}}}{\mathcal{V}m\pmb{v}} = \rho - \rho_{n}
\end{equation*}
而 $\rho_{n}$ 可以看作正常流体（normal fluid）的密度。由式 (3.7.13) 我们看到，零温下 $\rho_{s}=\rho$，尽管 $k=0$ 态中玻色子的数目小于 $N$。

如果玻色子带电并与电磁规范场 $A_{\mu}$ 耦合，那么具有零规范势的扭转凝聚体（$\pmb{A}=0$，$\mathrm{e}^{\mathrm{i}m\pmb{v}\cdot\pmb{x}}\varphi_{0}$）与具有非零规范势的无扭转凝聚体（$\pmb{A}=m\pmb{v}$，$\varphi_{0}$）是规范等价的。在这种情形下，式 (3.7.12) 可以解释为：通过打开一个常规范势 $\pmb{A}=m\pmb{v}$ 来产生有限动量 $\pmb{P}_{\pmb{v}}$。由于伽利略不变性（Galileo invariance），动量密度 $\pmb{P}_{\pmb{v}}/\mathcal{V}$ 正比于流密度 $\pmb{j}=\pmb{P}_{\pmb{v}}/m\mathcal{V}$。于是式 (3.7.12) 可以改写为
\begin{equation*}
\pmb{j} = \frac{\rho_{s}}{m}\pmb{A}
\end{equation*}
这正是伦敦方程（London equation）。我们看到，有限温度下伦敦方程中的系数由 $\rho_{s}/m$ 给出，而式 (3.7.13) 使我们能够计算它对温度的依赖关系。

% ---- 书页 135 (p150) ----
\subsection*{问题 3.7.2}
非伽利略不变超流体。尽管上面的讨论集中于伽利略不变的超流体，但所得到的大多数结果同样适用于非伽利略不变的超流体。

\begin{enumerate}
\item 对具有常规范势 $\pmb{A}$ 的玻色子系统导出 XY 模型，假设凝聚体 $\varphi=$ 常数。求出低能模式的色散 $\epsilon_{\pmb{A}}(\pmb{k})$，并与上面得到的 $\epsilon_{\pmb{v}}(\pmb{k})$ 比较。
\item 重复上述计算，但现在在玻色子拉格朗日量中附加一项 $\frac{c}{2}|\dot{\varphi}|^{2}$，以破坏伽利略不变性。
\item 导出上述非伽利略不变系统的伦敦方程。（提示：你可以先导出零温下的伦敦方程，然后考虑有限 $T$ 下的激发如何修正流。注意动量为 $\pmb{k}$ 的激发对流的贡献由 $\frac{\partial}{\partial\pmb{A}}\epsilon_{\pmb{A}}(\pmb{k})$ 给出。）
\end{enumerate}

\subsection{隧穿与约瑟夫森效应}

我们已经看到，超流相中玻色子的格林函数在不同的维度下颇为不同。本节我们将讨论两个超流体或超导体之间的隧穿（tunneling）。我们将看到，隧穿实验使我们能够测量玻色子格林函数的性质。

考虑由 $H_{R}$ 与 $H_{L}$ 描写、并通过一个隧穿算符 $I$ 耦合起来的两个玻色子系统，使得
\begin{equation*}
H = H_{R} + H_{L} + \Gamma I + \Gamma I^{\dagger}
\end{equation*}
其中
\begin{equation*}
I = \varphi_{L}\varphi_{R}^{\dagger}
\end{equation*}
$\Gamma$ 描写隧穿振幅。在存在电磁规范场时，总哈密顿量需要改写为
\begin{equation*}
H = H_{R} + H_{L} + \Gamma\mathrm{e}^{-\mathrm{i}a}I + \Gamma\mathrm{e}^{+\mathrm{i}a}I^{\dagger}
\end{equation*}
其中 $a=\int_{L}^{R}\pmb{A}\,\mathrm{d}\pmb{x}$ 是矢量势沿隧穿结的积分。隧穿流算符由下式给出
\begin{equation*}
j_{T} = \frac{\partial}{\partial a}H = -\mathrm{i}\Gamma\left(I\mathrm{e}^{-\mathrm{i}a} - I^{\dagger}\mathrm{e}^{+\mathrm{i}a}\right)
\end{equation*}
若 $\langle j_{T}\rangle>0$，则电流从 L 流向 R。在 $A_{0}=0$ 规范中，两个系统之间的电压差 $V=V_{L}-V_{R}$ 可以通过令
\begin{equation*}
a(t) = \int_{L}^{R}\pmb{A}\,\mathrm{d}\pmb{x} = -Vt
\end{equation*}
而包括进来。$A_{0}=0$ 规范的优点是，接通电压并不影响 $H_{R,L}$。

% ---- 书页 136 (p151) ----
有了上述设置，我们就可以利用线性响应理论（见式 (2.2.3)）写出隧穿流的如下表达式：
\begin{align*}
\langle j_{T}\rangle(t) ={}& -\mathrm{i}\Gamma\int^{t}\mathrm{d}t'\ \Big\langle\big[j_{T}(t), \big(\mathrm{e}^{-\mathrm{i}a(t')}I(t') + \text{h.c.}\big)\big]\Big\rangle + \langle j_{T}\rangle_{0}(t)\\
={}& -\Gamma^{2}\int^{t}\mathrm{d}t'\ \mathrm{e}^{-\mathrm{i}a(t)}\Big(\mathrm{e}^{\mathrm{i}a(t')}\big\langle[I(t), I^{\dagger}(t')]\big\rangle + \mathrm{e}^{-\mathrm{i}a(t')}\big\langle[I(t), I(t')]\big\rangle + \text{h.c.}\Big)\\
& + \langle j_{T}\rangle_{0}(t)
\end{align*}
其中 $\langle j_{T}\rangle_{0}(t)$ 是在没有隧穿项时的平均。

如果两个玻色子系统都处于超流相或超导相，那么主要贡献来自 $\langle j_{T}\rangle_{0}(t)$：
\begin{equation*}
\langle j_{T}\rangle(t) = 2\Gamma|\varphi_{R}\varphi_{L}|\sin(\theta - a(t))
\end{equation*}
其中 $\theta$ 是两个凝聚体 $\varphi_{R}$ 与 $\varphi_{L}$ 之间的相位差。我们看到，即使在零电压 $a(t)=0$ 时，只要 $\theta\neq0$，隧穿流就可以不为零：
\begin{equation*}
\langle j_{T}\rangle = I_{c}\sin(\theta), \qquad I_{c} = 2\Gamma|\varphi_{R}\varphi_{L}|
\end{equation*}
其中 $I_{c}$ 是最大隧穿流（称为临界隧穿流）。

如果其中一个玻色子系统处于正常相（包括代数衰减相），则 $\langle j_{T}\rangle_{0}(t)=0$ 且 $\langle[I(t), I(t')]\rangle=0$。我们有
\begin{equation*}
\langle j_{T}\rangle(t) = -\Gamma^{2}\int^{t}\mathrm{d}t'\left(\mathrm{e}^{-\mathrm{i}a(t)+\mathrm{i}a(t')}\big\langle[I(t), I^{\dagger}(t')]\big\rangle + \text{h.c.}\right) \tag{3.7.15}
\end{equation*}
由于 $\big\langle[I(t), I^{\dagger}(t')]\big\rangle$ 是没有隧穿时的关联，它可以表示成结两侧格林函数的乘积。例如，
\begin{equation*}
\big\langle I(t)I^{\dagger}(t')\big\rangle = \big\langle\varphi_{L}(t)\varphi_{L}^{\dagger}(t')\big\rangle\big\langle\varphi_{R}^{\dagger}(t)\varphi_{R}(t')\big\rangle
\end{equation*}

在 $1+1$ 维中，编时格林函数具有如下的代数衰减：
\begin{equation*}
\big\langle\varphi_{L,R}(t)\varphi_{L,R}^{\dagger}(t')\big\rangle \sim \mathrm{e}^{-\mathrm{i}\eta_{L,R}\pi/2}\,|t-t'|^{-\eta_{L,R}}.
\end{equation*}

% ---- 书页 137 (p152) ----
我们发现隧穿流具有如下形式\footnote{我们用到了如下事实
\begin{gather*}
\big\langle\varphi_{L,R}(t)\varphi_{L,R}^{\dagger}(t')\big\rangle = \big\langle\varphi_{L,R}^{\dagger}(t)\varphi_{L,R}(t')\big\rangle,\\
\big\langle\varphi_{L,R}^{\dagger}(t')\varphi_{L,R}(t)\big\rangle = \big\langle\varphi_{L,R}^{\dagger}(t)\varphi_{L,R}(t')\big\rangle^{*}
\end{gather*}
来证明
\begin{align*}
\big\langle[I(t), I^{\dagger}(t')]\big\rangle ={}& \big\langle\varphi_{R}^{\dagger}(t)\varphi_{R}(t')\big\rangle\big\langle\varphi_{L}(t)\varphi_{L}^{\dagger}(t')\big\rangle - \big\langle\varphi_{R}(t')\varphi_{R}^{\dagger}(t)\big\rangle\big\langle\varphi_{L}^{\dagger}(t')\varphi_{L}(t)\big\rangle\\
={}& \mathrm{e}^{-\mathrm{i}(\eta_{R}+\eta_{L})\pi/2}\left(\frac{l}{v|t-t'|}\right)^{\eta_{R}+\eta_{L}} - \text{h.c.}\\
={}& -2\mathrm{i}\sin\left(\frac{(\eta_{R}+\eta_{L})\pi}{2}\right)\left(\frac{l}{v|t-t'|}\right)^{\eta_{R}+\eta_{L}}
\end{align*}}
\begin{align*}
\langle j_{T}\rangle ={}& 2\sin\left(\frac{(\eta_{R}+\eta_{L})\pi}{2}\right)\left(\frac{l}{v}\right)^{\eta_{R}+\eta_{L}}\\
&\times \Gamma^{2}\int^{t}\mathrm{d}t'\ \left(\mathrm{i}\,\mathrm{e}^{\mathrm{i}V(t-t')}(t-t')^{-\eta_{R}-\eta_{L}}\frac{1}{|t-t'|^{\eta_{R}+\eta_{L}}} + \text{h.c.}\right)
\end{align*}
它导致一条非线性的 $I$--$V$ 曲线\footnote{我们用到了如下事实
\begin{align*}
\int_{-\infty}^{0}\mathrm{d}t\ \mathrm{e}^{\mathrm{i}Vt}|t|^{-a} ={}& \Big|_{t\to-\mathrm{i}\tau\,\mathrm{sgn}(V)}\ -\mathrm{i}\,\mathrm{sgn}(V)\,\mathrm{e}^{\mathrm{i}a\pi\,\mathrm{sgn}(V)/2}\int_{0}^{+\infty}\mathrm{d}\tau\ \mathrm{e}^{-|V|\tau}\tau^{-a}\\
={}& -\mathrm{i}\,\mathrm{sgn}(V)\,\mathrm{e}^{\mathrm{i}a\pi\,\mathrm{sgn}(V)/2}\,|V|^{a-1}\Gamma(1-a)
\end{align*}
以及 $\Gamma(a)\Gamma(1-a) = \frac{\pi}{\sin(\pi a)}$。}
\begin{equation*}
\langle j_{T}\rangle = 2\left(\frac{l}{v}\right)^{\eta_{R}+\eta_{L}}\Gamma^{2}\frac{\pi}{\Gamma(\eta_{R}+\eta_{L})}\,|V|^{\eta_{R}+\eta_{L}-1}\,\mathrm{sgn}(V)
\end{equation*}
我们看到，代数衰减的指数 $\eta_{R}+\eta_{L}$ 可以在隧穿实验中测量。

\subsection*{问题 3.7.3}
求有限温度下隧穿 $I$--$V$ 曲线的表达式，并证明 $V=0$ 处的微分电导具有温度依赖关系 $\mathrm{d}I/\mathrm{d}V\propto T^{\eta_{R}+\eta_{L}-2}$。

\subsection{Anderson--Higgs 机制}
\begin{keypoints}
\item Anderson--Higgs 机制把无能隙的 Nambu--Goldstone 模与无能隙的规范模结合成一个具有有限能隙的模。
\end{keypoints}

在上面的讨论中，我们把电磁规范场 $A_{\mu}$ 当作一个背景场（background field），它是固定的，不随电荷

% ---- 书页 138 (p153) ----
与流的变化而响应。本节中，我们将把 $A_{\mu}$ 当作一个具有自身涨落的动力学场来处理。

玻色子与规范场的量子理论可以通过路径积分来定义
\begin{equation*}
Z = \int\mathcal{D}\varphi\,\mathcal{D}A_{\mu}\ \mathrm{e}^{\mathrm{i}\int \mathrm{d}^{d}\pmb{x}\,\mathrm{d}t\, L(\varphi,A_{\mu}) + \frac{1}{8\pi e^{2}}(\pmb{E}^{2}-\pmb{B}^{2})}
\end{equation*}
其中我们已令光速 $c=1$。这里 $L(\varphi,A_{\mu})$ 是式 (3.7.3) 给出的加了规范场的玻色子拉格朗日量。本节我们只考虑由
\begin{equation*}
L = L(\varphi,A_{\mu}) + \frac{1}{8\pi e^{2}}(\pmb{E}^{2}-\pmb{B}^{2}) \tag{3.7.16}
\end{equation*}
描写的经典理论。量子规范理论将在第 6 章讨论。

由于规范不变性 (3.7.4)，如果 $(\varphi(\pmb{x},t), A_{\mu}(\pmb{x},t))$ 是运动方程的一个解，那么规范变换后的场 $(\tilde{\varphi}(\pmb{x},t), \tilde{A}_{\mu}(\pmb{x},t))$ 也满足运动方程。在规范理论中，我们不把 $(\varphi(\pmb{x},t), A_{\mu}(\pmb{x},t))$ 与 $(\tilde{\varphi}(\pmb{x},t), \tilde{A}_{\mu}(\pmb{x},t))$ 视为两个不同的运动，而是把它们视为同一个运动。换言之，场 $(\varphi(\pmb{x},t), A_{\mu}(\pmb{x},t))$ 是物理运动的多对一的标记。两组规范等价的场 $(\varphi(\pmb{x},t), A_{\mu}(\pmb{x},t))$ 与 $(\tilde{\varphi}(\pmb{x},t), \tilde{A}_{\mu}(\pmb{x},t))$ 是描写同一运动的两个标记。

在对称性破缺相中，我们有 $|\varphi(\pmb{x},t)|\neq0$。我们可以通过一个规范变换使 $\varphi(\pmb{x},t)$ 成为实数，即 $\varphi(\pmb{x},t)=|\varphi(\pmb{x},t)|=\phi(\pmb{x},t)$。这一手续称为固定规范（fixing the gauge）或选取规范（choosing the gauge）。固定规范之后，不同的 $(\phi, A_{\mu})$ 对将描写物理上不同的运动。在 $\varphi(\pmb{x},t)$ 取实值的规范中，我们有
\begin{equation*}
L = -A_{0}\phi^{2} - \frac{1}{2m}(\partial_{i}\phi)^{2} - \frac{\phi^{2}}{2m}(A_{i})^{2} - V(\phi) + \frac{1}{8\pi e^{2}}(\pmb{E}^{2}-\pmb{B}^{2})
\end{equation*}
把 $\phi=\varphi_{0}+\delta\phi$ 的小涨落积掉之后，我们得到下列低能有效理论：
\begin{equation*}
L_{\mathrm{eff}} = \frac{1}{2V_{0}}(A_{0})^{2} - \frac{\rho}{2m}(A_{i})^{2} + \frac{1}{8\pi e^{2}}(\pmb{E}^{2}-\pmb{B}^{2}) \tag{3.7.17}
\end{equation*}
上式也可以从加了规范场的 XY 模型 (3.7.10) 出发、令 $\theta=0$ 而得到。注意到
\begin{equation*}
\frac{1}{8\pi e^{2}}\pmb{E}^{2} = \frac{1}{8\pi e^{2}}(\partial_{0}A_{i})^{2} + \frac{1}{8\pi e^{2}}(\partial_{i}A_{0})^{2} - \frac{1}{4\pi e^{2}}\partial_{i}A_{0}\partial_{0}A_{i}
\end{equation*}
而 $A_{0}$ 不含时间导数项。因此 $A_{0}$ 不是动力学变量，我们把它积掉，得到
\begin{equation*}
L_{\mathrm{eff}} = \partial_{0}A_{j}\left(\frac{1}{8\pi e^{2}}\delta_{ij} + \frac{1}{2}\,\frac{V_{0}}{(4\pi)^{2}e^{4}}\,\partial_{j}\partial_{i}\right)\partial_{0}A_{i} - \frac{\rho}{2m}A_{i}^{2} - \frac{1}{8\pi e^{2}}\pmb{B}^{2}
\end{equation*}

% ---- 书页 139 (p154) ----
引入横分量与纵分量如下：
\begin{equation*}
A_{i} = \hat{k}_{i}A^{\parallel} + \hat{n}_{a}A_{a}^{\perp} \tag{3.7.18}
\end{equation*}
其中 $\hat{k}$、$\hat{n}_{1}$ 与 $\hat{n}_{2}$ 构成一组局域正交基，我们便可把 $L_{\mathrm{eff}}$ 改写为
\begin{align*}
L_{\mathrm{eff}} ={}& \frac{1}{8\pi e^{2}}\left((\partial_{0}A_{a}^{\perp})^{2} - (\partial_{i}A_{a}^{\perp})^{2}\right) - \frac{\rho}{2m}(A_{a}^{\perp})^{2}\\
&+ \frac{1}{8\pi e^{2}}\,\partial_{0}A^{\parallel}\left(1 + \frac{V_{0}}{4\pi e^{2}}(\partial_{i})^{2}\right)\partial_{0}A^{\parallel} - \frac{\rho}{2m}(A^{\parallel})^{2}.
\end{align*}
由运动方程我们发现，$L_{\mathrm{eff}}$ 所描写的三个模式都具有相同的能隙 $\Delta = e\sqrt{4\pi\rho/m}$。不存在无能隙模式。无能隙规范模式与无能隙 Nambu--Goldstone 模式之间的耦合赋予这些模式有限的能隙。这一现象称为 Anderson--Higgs 机制（Anderson, 1963; Higgs, 1964）。我们还可以把 $\partial_{i}^{2}\partial_{0}^{2}$ 项中的 $\partial_{0}$ 换成 $-\mathrm{i}\Delta$，从而进一步简化 $L_{\mathrm{eff}}$：
\begin{align*}
L_{\mathrm{eff}} ={}& \frac{1}{8\pi e^{2}}\left((\partial_{0}A_{a}^{\perp})^{2} - (\partial_{i}A_{a}^{\perp})^{2}\right) - \frac{\rho}{2m}(A_{a}^{\perp})^{2}\\
&+ \frac{1}{8\pi e^{2}}\left((\partial_{0}A^{\parallel})^{2} - v^{2}(\partial_{i}A^{\parallel})^{2}\right) - \frac{\rho}{2m}(A^{\parallel})^{2} \tag{3.7.19}
\end{align*}
（注意 $v^{2}=V_{0}\rho/m$ 是 XY 模型的速度。）从 $L_{\mathrm{eff}}$ 出发，我们可以研究带电超流体的所有经典电磁性质。

\subsection*{问题 3.7.4}
由拉格朗日量 $L=\frac{1}{8\pi e^{2}}(c^{-1}\pmb{E}^{2}-c\pmb{B}^{2})$ 导出 Maxwell 方程。证明规范涨落是无能隙的，且 $c$ 就是光速。

\section{热势的微扰计算}

\subsection{微扰与费曼规则}
\begin{keypoints}
\item 相互作用玻色子系统的费曼图与费曼规则。
\end{keypoints}

本节我们将讨论如何系统地计算相互作用玻色子系统的热势。为简单起见，我们假设玻色子场是实的。一个复玻色子场可以视为由两个分量组成，其中每个分量都是一个实玻色子场。下面的讨论可以容易地推广到多分量玻色子场。我们从有限温度下的虚时路径积分出发。计算热势的第一步是找到一个经典解 $\varphi_{c}$，并把作用量 $S$ 绕它

% ---- 书页 140 (p155) ----
展开如下：
\begin{equation*}
S(\varphi) = S(\varphi_{c}) + \int_{0}^{\beta}\mathrm{d}x^{\mu}\mathrm{d}y^{\mu}\ \frac{1}{2}\delta\varphi(x^{\mu})\mathcal{K}^{\beta}(x^{\mu}-y^{\mu})\delta\varphi(y^{\mu}) + \int_{0}^{\beta}\mathrm{d}x^{\mu}\,V(\delta\varphi)
\end{equation*}
其中 $V(\delta\varphi)=g_{3}\delta\varphi^{3}+g_{4}\delta\varphi^{4}+\dots$。如果我们忽略高阶项 $V(\delta\varphi)$，我们就得到一个自由系统，其热势 $\Omega_{0}$ 可以通过高斯积分算出，如 3.7.5 节所讨论。为计算相互作用系统的热势，我们注意到
\begin{align*}
Z = A\int\mathcal{D}\varphi\ \mathrm{e}^{-S} ={}& \mathrm{e}^{-S(\varphi_{c})}\mathrm{e}^{-\beta\Omega_{0}}\frac{\int\mathcal{D}\delta\varphi\ \mathrm{e}^{-S_{0}-\int \mathrm{d}x^{\mu}V(\delta\varphi)}}{\int\mathcal{D}\delta\varphi\ \mathrm{e}^{-S_{0}}}\\
={}& \mathrm{e}^{-S(\varphi_{c})}\mathrm{e}^{-\beta\Omega_{0}}\sum_{n=0}\frac{(-)^{n}}{n!}\left\langle\left(\int \mathrm{d}x^{\mu}\,V(\delta\varphi)\right)^{n}\right\rangle_{0}
\end{align*}
其中 $\langle\dots\rangle_{0}$ 表示权重为 $\mathrm{e}^{-S_{0}}$ 的平均，而 $S_{0}$ 是作用量的二次部分，即 $S_{0}\allowbreak=\int_{0}^{\beta}\mathrm{d}x^{\mu}\mathrm{d}y^{\mu}\allowbreak\ \frac{1}{2}\delta\varphi(x^{\mu})\mathcal{K}^{\beta}(x^{\mu}-y^{\mu})\delta\varphi(y^{\mu})$。如果存在若干个经典解（或驻定路径）$\varphi_{c}^{(a)}$，那么我们应当把它们全部包括进来：
\begin{equation*}
Z = \sum_{a}\mathrm{e}^{-S(\varphi_{c}^{(a)})}\mathrm{e}^{-\beta\Omega_{0}^{(a)}}\frac{\int\mathcal{D}\delta\varphi\ \mathrm{e}^{-S_{0}^{(a)}-\int \mathrm{d}x^{\mu}V^{(a)}(\delta\varphi)}}{\int\mathcal{D}\delta\varphi\ \mathrm{e}^{-S_{0}^{(a)}}}
\end{equation*}

我们看到，为计算总热势，我们需要计算多点关联 $\langle\prod_{i=1}^{n}\delta\varphi(i)\rangle_{0}$。对 $n=2$，它就是玻色子场的格林函数 $\langle\delta\varphi(1)\delta\varphi(2)\rangle_{0}=\mathcal{G}^{\beta}(1,2)$，其中 1 与 2 分别代表第一个与第二个玻色子场的坐标 $x_{1}^{\mu}$ 与 $x_{2}^{\mu}$。作为一个算符，$\mathcal{G}^{\beta}(1,2)$ 是 $\mathcal{K}^{\beta}(1,2)$ 的逆，即 $\int \mathcal{G}^{\beta}(1,2)\mathcal{K}^{\beta}(2,3)=\delta(x_{1}^{\mu}-x_{3}^{\mu})$。这里我们还使用了缩写 $\int f(1)\equiv\int \mathrm{d}^{d}x_{1}^{\mu}f(x_{1}^{\mu})$、$\int f(1)g(1,2)\equiv\int \mathrm{d}^{d}x_{1}^{\mu}\mathrm{d}^{d}x_{2}^{\mu}f(x_{1}^{\mu})g(x_{1}^{\mu},x_{2}^{\mu})$ 等。引入生成泛函
\begin{equation*}
Z[j] = \frac{\int\mathcal{D}\delta\varphi\ \mathrm{e}^{-S_{0}-\int \mathrm{d}x^{\mu}j\delta\varphi}}{\int\mathcal{D}\delta\varphi\ \mathrm{e}^{-S_{0}}} = \mathrm{e}^{\frac{1}{2}\int j(1)\mathcal{G}^{\beta}(1,2)j(2)}
\end{equation*}
我们发现（对偶数的 $n$）
\begin{align*}
\left\langle\prod_{i=1}^{n}\delta\varphi(i)\right\rangle_{0} ={}& \frac{\delta^{n}}{\delta j(1)\dots\delta j(n)}Z[j]\Big|_{j=0}\\
={}& \frac{1}{2^{n/2}(n/2)!}\ \frac{\delta^{n}}{\delta j(1)\dots\delta j(n)}\left(\int j\mathcal{G}^{\beta}j\right)^{n/2}\\
={}& \left(\mathcal{G}^{\beta}(1,2)\mathcal{G}^{\beta}(3,4)\dots\mathcal{G}^{\beta}(n-1,n)\right) + \text{所有不同的排列}
\end{align*}
上式是 Wick 定理的另一种形式。

% ---- 书页 141 (p156) ----
%FIG{3.16}: {两个顶点，代表 $\int g_{4}\delta\varphi^{4}(1)\int g_{4}\delta\varphi^{4}(2)$。}

利用 Wick 定理我们发现，例如 $\big\langle\int V(\delta\varphi)\int V(\delta\varphi)\big\rangle_{0}$ 包含许多项：
\begin{align*}
\left\langle\int V(\delta\varphi)\int V(\delta\varphi)\right\rangle_{0} ={}& \left\langle\int g_{4}\delta\varphi^{4}(1)\int g_{4}\delta\varphi^{4}(2)\right\rangle_{0} + \dots\\
={}& \left[4!g_{4}^{2}\int\left(\mathcal{G}^{\beta}(1,2)\right)^{4}\right] + \left[(6^{2}\times 2)g_{4}^{2}\int\left(\mathcal{G}^{\beta}(1,2)\right)^{2}\mathcal{G}^{\beta}(1,1)\mathcal{G}^{\beta}(2,2)\right]\\
&+ \left[3g_{4}\int\left(\mathcal{G}^{\beta}(1,1)\right)^{2}\right]\left[3g_{4}\int\left(\mathcal{G}^{\beta}(2,2)\right)^{2}\right] + \dots \tag{3.8.1}
\end{align*}
积分因子 $4!$、$6^{2}\times2$ 等来自 Wick 展开中恰好具有相同形状的不同项。例如，在一个更简单的情形中，Wick 展开 $\langle ABCD\rangle=\langle AB\rangle\langle CD\rangle+\langle AC\rangle\langle BD\rangle+\langle AD\rangle\langle BC\rangle$ 导致
\begin{align*}
\langle AABB\rangle ={}& \langle AA\rangle\langle BB\rangle + \langle AB\rangle\langle AB\rangle + \langle AB\rangle\langle AB\rangle\\
={}& \langle AA\rangle\langle BB\rangle + 2\langle AB\rangle\langle AB\rangle
\end{align*}
更方便的做法是利用费曼图，通过费曼规则直接得到 Wick 展开 (3.8.1) 的结果。我们不在一般的框架中描述费曼图与费曼规则，而是选择用式 (3.8.1) 给出的具体例子来解释它们。我觉得在实战中理解费曼图与费曼规则更容易。我们在这里要描述的费曼规则是实空间（real-space）费曼规则。动量空间费曼规则将在 5.4.2 节中描述。

为计算 $\big\langle\int g_{4}\delta\varphi^{4}(1)\int g_{4}\delta\varphi^{4}(2)\big\rangle_{0}$，我们从图 3.16 中代表 $\int g_{4}\delta\varphi^{4}(1)\int g_{4}\delta\varphi^{4}(2)$ 的两个顶点出发。为得到期望值 $\langle\dots\rangle_{0}$，我们需要把顶点的所有“腿”彼此连接起来。这就生成了图 3.17 中的图形。这些图形称为费曼图。为构造式 (3.8.1) 的结果，我们只需使用以下三条费曼规则。

(i) 费曼图中的每条线代表一个传播子 $\mathcal{G}^{\beta}(a,b)$，其中 $a$ 与 $b$ 是该线两端的两个顶点。

(ii) 每个顶点代表 $g_{4}\int$。

(iii) 连通图的值由规则 (i) 与 (ii) 给出，而非连通图

% ---- 书页 142 (p157) ----
%FIG{3.17}: {对应于式 (3.8.1) 中三项的三个费曼图。这里 (a) 与 (b) 是连通图，(c) 是包含两个连通图的非连通图。}

的值由连通子图的乘积给出。这样，我们从图 3.17(a)、(b) 与 (c) 所示的三个费曼图分别得到式 (3.8.1) 末行中的第一、第二与第三项。

嗯，并不完全如此。我们忽略了诸如 $4!$、$6^{2}\times2$ 等数值因子。这些数值因子是费曼规则中最困难的部分。它们从何而来？我们先考虑与图 3.17(a) 中图形相联系的因子 $4!$。这个因子来自这样一个事实：每个顶点有四条腿，而把顶点 1 的四条腿与顶点 2 的四条腿连接起来共有 $4!$ 种不同的方式。与图 3.17(b) 中图形相联系的因子 $6^{2}\times2$ 更复杂一些，但它仍然代表八条腿彼此连接的不同方式。首先，顶点 1 的两条腿彼此相连。从四条腿中选取两条有六种方式，这样我们就得到 $6^{2}\times2$ 中的因子 6。第二个因子 6 类似地来自顶点 2。把顶点 1 剩下的两条腿与顶点 2 剩下的两条腿连接起来有两种方式，于是我们得到最后的因子 2。第三项中的因子 3 对应于把一个顶点的四条腿分成两组、每组各有两条腿的不同分法。（注意这不同于从四条腿中选取两条腿的不同选法。）

让我们回到想要计算的热势。在得到所有的平均之后，相互作用对热势的修正可以这样求出：
\begin{equation*}
\delta\Omega = -T\ln\sum_{n=0}\frac{(-)^{n}}{n!}\left\langle\left(\int \mathrm{d}x^{\mu}\,V(\delta\varphi)\right)^{n}\right\rangle_{0}
\end{equation*}
有一个连链定理（linked-cluster theorem）可以简化上述计算。根据该定理，$\delta\Omega$ 由所有连通图之和给出：
\begin{equation*}
\delta\Omega = -T\sum_{n=0}\frac{(-)^{n}}{n!}\left\langle\left(\int \mathrm{d}x^{\mu}\,V(\delta\varphi)\right)^{n}\right\rangle_{0c}
\end{equation*}
其中 $\big\langle\big(\int \mathrm{d}x^{\mu}V(\delta\varphi)\big)^{n}\big\rangle_{0c}$ 是从 $\big\langle\big(\int \mathrm{d}x^{\mu}V(\delta\varphi)\big)^{n}\big\rangle_{0}$ 中去掉所有非连通图的贡献后得到的。

% ---- 书页 143 (p158) ----
\subsection{连链定理}
\begin{keypoints}
\item 只对连通的费曼图求和，就可以直接计算有效作用量。
\end{keypoints}

我们将用复制（replica）技巧来导出连链定理，这既是为了简洁，也是为了介绍这一有用的方法。复制方法的基本思想是，对整数 $n$，通过把系统复制 $n$ 份来计算 $Z^{n}$：
\begin{equation*}
\left(\frac{Z}{Z_{0}}\right)^{n} = \frac{\int\mathcal{D}\delta\varphi_{\alpha}\ \mathrm{e}^{-\sum_{\alpha=1}^{n}\left[S_{0}(\delta\varphi_{\alpha}) + \int \mathrm{d}x^{\mu}V(\delta\varphi_{\alpha})\right]}}{\int\mathcal{D}\delta\varphi_{\alpha}\ \mathrm{e}^{-\sum_{\alpha=1}^{n}S_{0}(\delta\varphi_{\alpha})}}
\end{equation*}
现在 $(Z/Z_{0})^{n}$ 可以通过微扰展开来计算。在每一个费曼图中，每条传播子都携带一个指标 $\alpha$，并且进入与离开同一个相互作用顶点的所有传播子都具有相同的指标。当我们把 $\alpha$ 从 1 求和到 $n$ 时，显然每个连通图都正比于 $n$，而每个非连通图正比于 $n^{N_{c}}$，其中 $N_{c}$ 是该非连通图中连通图的数目。因此，
\begin{align*}
\left(\frac{Z}{Z_{0}}\right)^{n} ={}& \mathrm{e}^{n\ln\frac{Z}{Z_{0}}} = 1 + n\ln\frac{Z}{Z_{0}} + \sum_{m=2}^{\infty}\frac{(n\frac{Z}{Z_{0}})^{m}}{m!}\\
={}& 1 + n\sum(\text{所有连通图}) + O(n^{2})
\end{align*}
这样我们就证明了连链定理。

\subsection*{问题 3.8.1}
考虑一个非简谐振子 $L=\frac{1}{2}m\dot{x}^{2}-\frac{1}{2}m\omega_{0}^{2}x^{2}-gx^{3}$。计算该系统精确到 $g^{2}$ 阶的有限温度自由能。（提示：在 $\tau$ 空间中做计算可能更容易。）

注意，在微扰论的框架内，上述非简谐振子是一个行为良好的系统。不稳定性来自反弹（bounce）。证明在小 $g$ 极限下，基态的衰变率具有 $\omega_{0}\mathrm{e}^{-\#m^{3}\omega_{0}^{5}/g^{2}}$ 的量级。这样的项不可能出现在围绕基态 $x=0$ 的微扰计算中。然而，反弹作为另一条驻定路径，确实对自由能有贡献。证明反弹的贡献以正比于 $\mathrm{e}^{-\#m^{3}\omega_{0}^{5}/g^{2}}$ 的非微扰项出现。

% ==================================================================
% 新增术语（ch3_g 新增，供术语表追加）：
%   two-fluid picture 二流体图像
%   superfluid density / normal fluid (density) 超流密度 / 正常流体（密度）
%   London equation 伦敦方程
%   Galileo invariance / Galileo non-invariant 伽利略不变性 / 非伽利略不变的
%   twist 扭转
%   tunneling / tunneling junction / tunneling amplitude 隧穿 / 隧穿结 / 隧穿振幅
%   tunneling current / critical tunneling current 隧穿流 / 临界隧穿流
%   algebraic-decay phase 代数衰减相
%   time-ordered Green's function 编时格林函数
%   I–V curve / differential conductance I–V 曲线 / 微分电导
%   background field / dynamical field 背景场 / 动力学场
%   fixing the gauge / choosing the gauge 固定规范 / 选取规范
%   transverse / longitudinal component 横分量 / 纵分量
%   Feynman diagram / Feynman rules 费曼图 / 费曼规则
%   real-space / momentum-space Feynman rules 实空间 / 动量空间费曼规则
%   vertex / leg（顶点的）腿  顶点
%   connected / disconnected diagram 连通图 / 非连通图
%   linked-cluster theorem 连链定理（术语表已有）
%   replica technique 复制技巧
%   anharmonic oscillator 非简谐振子
%   Maxwell equation Maxwell 方程
%   thermal potential 热势（术语表已有）
% ==================================================================
